    %-----------------------------------------------------------------------------------------------------------------------------------------------%
    % The MIT License (MIT)
    %
    % Copyright (c) 2021 Jitin Nair
    %
    % Permission is hereby granted, free of charge, to any person obtaining a copy
    % of this software and associated documentation files (the "Software"), to deal
    % in the Software without restriction, including without limitation the rights
    % to use, copy, modify, merge, publish, distribute, sublicense, and/or sell
    % copies of the Software, and to permit persons to whom the Software is
    % furnished to do so, subject to the following conditions:
    %
    % THE SOFTWARE IS PROVIDED "AS IS", WITHOUT WARRANTY OF ANY KIND, EXPRESS OR
    % IMPLIED, INCLUDING BUT NOT LIMITED TO THE WARRANTIES OF MERCHANTABILITY,
    % FITNESS FOR A PARTICULAR PURPOSE AND NONINFRINGEMENT. IN NO EVENT SHALL THE
    % AUTHORS OR COPYRIGHT HOLDERS BE LIABLE FOR ANY CLAIM, DAMAGES OR OTHER
    % LIABILITY, WHETHER IN AN ACTION OF CONTRACT, TORT OR OTHERWISE, ARISING FROM,
    % OUT OF OR IN CONNECTION WITH THE SOFTWARE OR THE USE OR OTHER DEALINGS IN
    % THE SOFTWARE.
    %
    %-----------------------------------------------------------------------------------------------------------------------------------------------%

    %----------------------------------------------------------------------------------------
    % DOCUMENT DEFINITION
    %----------------------------------------------------------------------------------------

    \documentclass[a4paper,10.999999pt]{article}

    %----------------------------------------------------------------------------------------
    % PACKAGES
    %----------------------------------------------------------------------------------------
    \usepackage[utf8]{inputenc}
    \usepackage{url}
    \usepackage{parskip}

    % other packages for formatting
    \RequirePackage{color}
    \RequirePackage{graphicx}
    \usepackage[usenames,dvipsnames]{xcolor}
    \usepackage[scale=0.85, top=0.85cm, bottom=0.85cm, left=0.85cm, right=0.85cm]{geometry}

    % tabularx environment
    \usepackage{tabularx}

    % for lists within experience section
    \usepackage{enumitem}
    \usepackage{setspace}

    % centered version of 'X' col. type
    \newcolumntype{C}{>{\centering\arraybackslash}X}

    % to prevent spillover of tabular into next pages
    \usepackage{supertabular}
    \usepackage{tabularx}
    \newlength{\fullcollw}
    \setlength{\fullcollw}{0.47\textwidth}

    % custom \section
    \usepackage{titlesec}
    \usepackage{multicol}
    \usepackage{multirow}

    % section formatting - BALANCED SPACING
    \titleformat{\section}{\large\scshape\raggedright}{}{0em}{}[\titlerule]
    \titlespacing{\section}{0pt}{5pt}{3pt}

    % global spacing tweaks - BALANCED
    \setlength{\parskip}{1pt}
    \setlength{\parindent}{0pt}
    \setlist[itemize]{nosep, topsep=1pt, partopsep=0pt, leftmargin=1em, itemsep=1pt}
    \linespread{0.962}\selectfont

    \usepackage[unicode, draft=false]{hyperref}
    \definecolor{linkcolour}{rgb}{0,0.2,0.6}
    \hypersetup{colorlinks,breaklinks,urlcolor=linkcolour,linkcolor=linkcolour}

    % job listing environments - BALANCED VERTICAL SPACE
    % Plain \hfill-aligned lines instead of tabularx: ATS parsers frequently
    % read table cells out of visual order (e.g. pull the date before the
    % title). A single line with \hfill renders identically but extracts as
    % ordinary left-to-right text.
    \newenvironment{jobshort}[2]
        {
        \noindent\textbf{#1}\hfill{#2}\par
        }
        {}

    \newenvironment{joblong}[2]
        {
        \noindent\textbf{#1}\hfill{#2}\par
        \vspace{-3.2pt}
        \begin{minipage}[t]{\linewidth}
        \begin{itemize}[nosep, after=\strut, leftmargin=0.5em, itemsep=1pt, label=\raisebox{0.2ex}{\tiny$\circ$}]
        }
        {
        \end{itemize}
        \end{minipage}
        \vspace{-3.2pt}
        }

    % Custom nested itemize style
    \newcommand{\subbulletitem}[1]{
      \item[\raisebox{0.20ex}{\small$-$}] #1
    }

    %----------------------------------------------------------------------------------------
    % BEGIN DOCUMENT
    %----------------------------------------------------------------------------------------
    \begin{document}

    \pagestyle{empty}

    %----------------------------------------------------------------------------------------
    % TITLE
    %----------------------------------------------------------------------------------------


    \begin{center}
    \LARGE{\textbf{Kelvi Manavadaria}} \\[4pt]
    \href{mailto:kelvimanavadaria@gmail.com}{kelvimanavadaria@gmail.com} \ $|$ \
    +91 7862036686 \ $|$ \
    \href{https://github.com/KelviManavadaria-17}{github.com/KelviManavadaria-17} \ $|$ \
    \href{https://www.linkedin.com/in/kelvi-manavadaria/}{linkedin.com/in/kelvi-manavadaria}
    \end{center}

    %----------------------------------------------------------------------------------------

    \section{Summary}
DevOps Engineer with 2+ years building CI/CD pipelines, GitOps automation, and Kubernetes-based deployments on AWS-first, multi-cloud infrastructure. Skilled in Terraform, container orchestration, and cloud IAM/security, with a track record of shipping highly available systems.

    % SKILLS SECTION
    %----------------------------------------------------------------------------------------

    \section{Skills}
    \begin{itemize}[leftmargin=1.5em, itemsep=2pt, label=\raisebox{0.25ex}{\tiny$\circ$}]
    \item \textbf{CI/CD \& GitOps:} Jenkins, GitLab CI, ArgoCD, GitHub Actions, AWS CodePipeline, AWS CodeBuild, GitOps
    \item \textbf{Containerization \& Orchestration:} Kubernetes, Docker, Helm
    \item \textbf{Cloud:} AWS (EC2, ECR, IAM, RDS, DynamoDB, Lambda, Glue, Athena, API Gateway, Cognito, EventBridge, CloudTrail), Azure (Azure Monitor)
    \item \textbf{Infrastructure as Code (IaC) \& Configuration Management:} Terraform, Ansible
    \item \textbf{Observability:} Prometheus, Grafana, New Relic, CloudWatch, Elasticsearch (ELK Stack)
    \item \textbf{SRE Practices:} SLA/SLO/SLI, Error Budgets, Incident Management, Root Cause Analysis (RCA), 24x7 On-Call, MTTR Reduction
    \item \textbf{Systems:} Linux, Networking, Distributed Systems
    \item \textbf{Architecture:} Microservices, High Availability, Fault Tolerance, Disaster Recovery
    \item \textbf{Languages, Scripting \& Delivery:} Python, Java, Bash, Agile/Scrum
    \end{itemize}

    %----------------------------------------------------------------------------------------
    % WORK EXPERIENCE
    %----------------------------------------------------------------------------------------

    \section{Work Experience}

    \begin{joblong}{DevOps Engineer}{Dec 2023 -- Present}
    Intuitive.Cloud
    \item \textbf{Outsourced as Platform Engineer at Thomson Reuters}
    \begin{itemize}
    \item Deployed indirect tax calculation microservices on Kubernetes clusters using Helm, enabling scalable, fault-tolerant, and highly available systems across environments via CI/CD pipelines, and architected a multi-layer observability stack for proactive alerting.
    \item Automated workflows following GitOps practices using GitHub Actions for building custom Docker images, and built authentication for secure image access to Amazon ECR using AWS Cognito.
    \end{itemize}
    \vspace{5pt}

    \item \textbf{Site Reliability Engineer (SRE) in Cloud Managed Services Team}
    \begin{itemize}
    \item Applied SRE practices including 24x7 on-call production support, incident management, root cause analysis (RCA), and proactive reliability improvements, ensuring adherence to SLA/SLO targets and error budgets, improving system stability in a high-availability distributed environment.
    \item Designed a monitoring solution leveraging New Relic for enhanced observability across 30+ accounts.
    \item Automated alerting using CloudWatch, EventBridge, CloudTrail, Lambda, aligned with SLI/SLO and proactive incident management.
    \item Built hybrid monitoring with Azure Monitor, ensuring high availability and supporting business continuity across multi-cloud environments.
    \item Processed 10+ TB/day structured and unstructured data using AWS Glue and ETL pipelines for real-time analytics.
    \item Built querying system using AWS Athena, AWS Lambda, API Gateway, reducing MTTR and improving debugging efficiency.
    \end{itemize}
    \end{joblong}

    \begin{joblong}{Cloud Engineer Intern}{May 2023 -- Jun 2023}
    Mactores Innovation
    \item \textbf{High-Availability CI/CD Pipeline in AWS}
    \begin{itemize}
    \item Implemented a resilient CI/CD pipeline using AWS CodePipeline, CodeBuild, and GitHub Actions, ensuring 99.95\% uptime.
    \end{itemize}
    \end{joblong}

    \begin{joblong}{DevOps Intern}{May 2022 -- May 2023}
    Motorola Solutions
    \item \textbf{Scalable Logging and Monitoring with ELK Stack}
    \begin{itemize}
    \item Built centralized logging and monitoring system using ELK Stack and Docker, improving observability and debugging efficiency.
    \end{itemize}
    \end{joblong}

    %----------------------------------------------------------------------------------------
    % PROJECTS SECTION
    %----------------------------------------------------------------------------------------

    \section{Projects and Blogs}

    \begin{itemize}[leftmargin=1.5em, itemsep=2pt, label=\raisebox{0.25ex}{\tiny$\circ$}]
    \item \textbf{Other Projects} \href{https://github.com/KelviManavadaria-17}{[Github]}
    \begin{itemize}[nosep, leftmargin=1.5em, itemsep=1pt, label=\raisebox{0.25ex}{\small$-$}]
    \item Configured a CI/CD pipeline integrating Jenkins, GitHub Actions, and Argo CD for automated Kubernetes deployments.
    \end{itemize}
    \item \textbf{Blogs} \href{https://medium.com/@KelviManavadaria}{[Medium]}
    \end{itemize}

    %----------------------------------------------------------------------------------------
    % EDUCATION SECTION
    %----------------------------------------------------------------------------------------

    \section{Education}

    \noindent\textbf{B. Tech in Computer Engineering}\hfill{2020 -- 2024}\par
    \noindent Charotar University of Science and Technology\hfill{\textbf{GPA: 9.55}}\par

    \vspace{3pt}

    %----------------------------------------------------------------------------------------
    % CERTIFICATIONS SECTION
    %----------------------------------------------------------------------------------------

    \section{Certifications}

    \begin{itemize}[leftmargin=1.5em, itemsep=2pt, label=\raisebox{0.25ex}{\tiny$\circ$}]
    \item AWS: Solutions Architect – Associate (Expired Jun 2026) | SysOps Administrator – Associate (Valid thru Apr 2027) | Cloud Practitioner (Valid thru Apr 2027)
    \item Certified Kubernetes Administrator (CKA) — Valid thru Nov 2026
    \item Multi-Cloud: Google Cloud Digital Leader (Valid thru Feb 2027) | Microsoft Certified: Azure Fundamentals (Issued Dec 2022)
    \item HashiCorp Certified: Terraform Associate (003) (Expired Oct 2025)
    \item Additional Certifications (incl. AWS Partner: Sales \& Technical Accredited): https://www.credly.com/users/kelvi-manavadaria
    \end{itemize}

    \vspace{1.5pt}

    \end{document}
